\section{State of the Art}

\subsection{Delta Storage Systems}

Several programs have been developed for adding versioning and incremental updates 
to the write-only Parquet storage files. 

The program closest to meeting the needs of the OpenINTEL storage system is 
Apache Hudi\footnote{\url{https://hudi.apache.org/}}. Hudi is a storage management 
system built on top of database formats like Parquet. Hudi adds 
log files onto these files, which describe changes and can be viewed 
as deltas to the original Parquet file. However, it is limited by its retrieval and 
storage capabilities. Hudi was created as an efficient way to apply edits to a 
write-only database. Because of this, it only has functions to retrieve the current status 
which applies all delta logs at once. Hudi lacks the feature of constructing the table for 
any given snapshot. Furthermore, a built-in feature that cannot be removed from Hudi is its 
Compaction feature. Hudi will regularly update the base snapshot which merges the first few 
log files into the base and writes a new base file. This would destroy the history stored in 
the log files.

Both Apache Iceberg\footnote{\url{https://iceberg.apache.org/}} and Delta Lake\footnote{\url{https://delta.io/}} 
take a different approach to extending the Parquet storage format. Both systems create a time-travel 
system by keeping track of snapshots for every update. This is done by creating new Parquet 
files for every update to the base. A timestamp can then be added to the query to which these systems 
retrieve the relevant Parquet files. However, the important pitfall with these systems is that 
they are not capable of storing deltas; they only store snapshots. Unlike Hudi, these systems 
are not designed to progressively add snapshots together, but for full snapshots. These systems 
expect every measurement to be a full snapshot repeating all the data and not saving on storage 
efficiency at all. If these systems would be used by adding only the changed rows to each day, 
the systems are not capable of combining the different snapshots together, this would have to be 
done manually where each Parquet for each day would be retrieved and combined.

Furthermore, none of these current systems are able to handle formats with partially 
volatile columns that update every measurement.

\subsection{DNS Measurement Storage Systems}

Next to the OpenINTEL project which this research will build on, several DNS 
history providers exist. The most prominent one is the ActiveDNS 
project\footnote{\url{https://activednsproject.org/}}. Similar to OpenINTEL, ActiveDNS stores 
daily DNS records. The ActiveDNS project collects data from multiple vantage points; 
this results in possible duplicate data which is deduplicated in the storage module 
of the system~\cite{DBLP:conf/raid/KountourasKLCND16}. However, like OpenINTEL, 
ActiveDNS does not have any deduplication of data between separate days. 

Furthermore, commercial providers exist such as the Farsight 
DNSDB\footnote{\url{https://www.domaintools.com/}}, WhoisXMLAPI's DNS 
Database\footnote{\url{https://dns-history.whoisxmlapi.com/database}}, and SpamHaus's
passive DNS database\footnote{\url{https://www.spamhaus.com/data-access/passive-dns-api/}}. 
None of these providers provide any public documentation on their storage system, thus 
it cannot be concluded if any type of delta-compression is used to efficiently store their data.

\subsection{Delta Strategies}

Delta storage is used in many video codecs used to efficiently store video data, such as the MPEG video encoding 
standards~\cite{isoiec11172}. The MPEG codecs store frames of a video as a combination 
of base images, forward predicted images, and bi-directionally predicted images. The 
bi-directionally predicted images allow for even smaller deltas as the similarities to both 
the previous and the next data point can be used. Under the hypothesis that DNS records 
do not change often, this could be beneficial for the OpenINTEL storage system in principle.
However, the OpenINTEL storage system does not have immediate access to all data. 
The system would not be able to store the current day's data without knowing the data of the 
next day. This would cause more temporary storage overhead and adds complexity to the 
storage system which is unfavorable.

Furthermore, Molfetas et al.~\cite{DBLP:conf/awc/MolfetasWZ14} tested several combinations of files and strategies 
with the goal of reducing the amount of dependencies on other snapshots while retaining storage efficiency. 
This work is similar in nature since it touches on the effects of using delta-compression, however it does not take into account 
the vast amount of the data being stored by the OpenINTEL project and the significant time loss of 
I/O operations on the hard disks for retrieval of the resulting data that this creates.

LoopDelta, a delta-based compression framework with backups in mind, does identify the extra I/O 
operations necessary for reconstruction of delta files as a bottleneck in delta-based 
storage~\cite{10.1145/3721485}. But, similar to the aforementioned video codecs, LoopDelta relies 
partially on backwards deltas which require the next day to be known in the OpenINTEL system 
for the compression optimization.
